\section{Linear endpoint descent with original-position certificates}
\label{sec:linear-braid-descent}

The endpoint kernel of report 34 removes a repeated whole-word scan from the
maintained braid recognizer. Previously, every elementary destabilization
ran cyclic free reduction, counted both endpoint generators, copied a cyclic
rotation, and potentially renumbered every surviving letter. A chain on
$m$ strands and $n$ letters could therefore use $O(nm)$ word operations.
The replacement retains the same sufficient endpoint criterion and the
same right-before-left preference. It does not find general Markov moves,
internal connected-sum cuts, or a minimum braid index.

\subsection{Representation and local updates}

Keep the original signed generators $w_i$, with predecessor and successor
arrays forming a cyclic list of live positions. For each original generator
index, maintain its live multiplicity and the xor of its live positions.
The xor is used only when the multiplicity is one, when it identifies that
unique position exactly. A live interval of strand labels is represented by
$\ell<h$: the current strand count is $h-\ell$, and an original signed
letter becomes
\[
 \operatorname{sgn}(w_i)(|w_i|-\ell).
\]
A left destabilization increments $\ell$; a right destabilization decrements
$h$. No surviving word is renumbered until output.

Initially queue every original position as a possible left endpoint of an
inverse seam. Remove a live adjacent inverse pair when found, then enqueue
only the predecessor of the newly exposed seam. After the queue drains,
the cyclic word has no inverse adjacency. If the current right endpoint
$h-1$ occurs once, remove it; otherwise test the left endpoint $\ell+1$.
Set the head to the successor of the removed letter, implementing precisely
the cyclic conjugation required to put that crossing last before Markov
destabilization. Enqueue the newly exposed seam and repeat. Stop at three
strands or when neither endpoint is a singleton.

Every inverse cancellation is a braid-group equality, including a conjugation
for a wrapped cyclic cancellation. An endpoint generator occurring exactly
once can be moved to the last position by conjugation and destabilized;
left renumbering merely identifies the surviving consecutive strand labels.
Thus every recorded move preserves closure isotopy. Reduction order can
change the chosen cyclic representative; comparison with the earlier
algorithm is therefore up to cyclic rotation.

\begin{proposition}[Amortized endpoint cost]
On an explicit braid with $n$ letters and $m$ strands, the linked kernel and
its original-position replay use $O(n+m)$ word operations and space. Their
conservative bit cost, including the certificate, is
$O((n+m)\log(n+m+2))$.
\end{proposition}
\begin{proof}
Each position is deleted at most once. A cancellation deletes two positions
and a destabilization deletes one; each deletion event adds at most one queue
entry. Including the initial $n$ entries gives at most $2n$ queue pops and
at most $n$ trace records. Each pop and deletion performs a bounded number
of array, counter and link operations. Every successful strand iteration
removes a letter, and there is at most one final unsuccessful iteration.
Initialization costs $O(n+m)$ and final output costs $O(n)$. Replay checks
only the supplied local adjacencies and singleton counts, with the same
initialization and deletion bounds. Indices and signed generator labels
have $O(\log(n+m+2))$ bits, giving the conservative bit bound.
\end{proof}

The maintained dispatcher selects the linked path for at least 64 strands
and 128 letters, provided $m\le n+1$. Otherwise it retains the simpler
version-1 implementation. A one-component closure necessarily satisfies
$m\le n+1$, since each adjacent transposition can reduce the number of
permutation cycles by at most one. On this domain the remaining branch has
bounded rank or bounded word length, so its whole-word passes also have a
linear bound with fixed constants. Sparse standalone link inputs retain
the earlier $O(nm)$ upper bound and avoid an allocation indexed by an
arbitrarily large strand count. These statements concern explicit input;
they are not complexity claims for a binary-compressed braid word.

\subsection{Replay and resource boundaries}

The new \texttt{singleton-markov-descent-v2} certificate records original
positions for every cancellation and destabilization, plus the final live
order and signed word. The verifier constructs its own lists and multiplicities;
it checks live distinct positions, directed adjacency, inverse generators,
endpoint identity and singleton count. It does not use the producer's xor
lookup or queue, and it does not search for a reducible endpoint. Exact
field, integer and trace-length checks reject malformed records. The queue
counter is a diagnostic, never a mathematical witness for an isotopy.
A valid prematurely stopped trace proves only its recorded equivalence.

Version-1 evidence remains replayable by the old verifier. Its historical
replay complexity is unchanged; the new linear replay bound applies to the
new certificates. Both new production and replay paths call the cooperative
cancellation callback during their traversals. Resource interruption raises
without returning a partial recognition result.

The complete free-product decision for the resulting at-most-three-strand
braid remains linear on explicit input, so this endpoint-recognized class
now has a linear symbolic path. The independent matrix backend retains its
quadratic bit-arithmetic term. A supplied PD still needs checked braid
provenance; neither this improvement nor a stalled endpoint search makes a
new general unknot decision.

\subsection{Validation and measurement contract}

Seven added test methods compare exhaustive signed words through three
strands, additional four-strand cases, and 120
seeded conjugated mixed stabilization chains with a separate explicit
whole-word algorithm. A deliberately slow list-based original-position
oracle checks the local trace independently of both linked implementations.
Additional tests cover both signs, both recognition backends, chains through
2048 strands, malformed and reused positions, false inverse cancellations,
input binding, small-path dispatch and resource interruption. Small
stabilizations agree with independently computed reduced Khovanov rank.
The complete maintained suite passes all 775 tests in 181.567 seconds with
Regina installed.

\path{fast/benchmark_braid_descent.py} pins the previous reducer and verifier
to commit \texttt{c4c23b61d}. It compares baseline, identical A/A control,
and the adaptive implementation in shuffled order, using one excluded
warmup and five measured rounds. Kernel measurements include certificate
generation and independent replay. Raw braid measurements also include
closure validation and symbolic recognition; fresh-PD measurements include
diagram construction and the ordinary recognition pipeline. Input construction
for the kernel is excluded; imports and JSON serialization are excluded in
all arms. Source hashes are checked before and after the experiment. Any
unknown result suppresses a timing ratio for that case.

The experiment contains 300 completed measured calls and 60 excluded warmups.
At 4096 strands, right-chain generation plus replay falls from 4300.024 ms
to 14.142 ms, with median paired ratio 305.29; the left chain falls from
5953.078 ms to 14.304 ms, with paired ratio 380.36. At 256 strands the
corresponding ratios are 19.10 and 24.27. Complete raw-braid recognition
plus replay at 2048 strands falls from 1042.500 ms to 7.582 ms, with paired
ratio 138.34. These are large separated effects, supported independently
by the deletion and queue bounds rather than extrapolated from timings.

\begin{center}\small
\begin{tabular}{@{}lrrrr@{}}
\toprule
Case & Before & A/A & Adaptive & Ratio \\
\midrule
right-16 & 0.127 & 0.128 & 0.128 & 0.99 \\
left-16 & 0.202 & 0.175 & 0.172 & 1.20 \\
right-64 & 1.248 & 1.187 & 1.196 & 1.05 \\
left-64 & 1.587 & 1.570 & 1.579 & 1.01 \\
right-128 & 4.280 & 4.216 & 4.259 & 0.99 \\
left-128 & 5.388 & 5.394 & 0.446 & 12.35 \\
right-256 & 15.470 & 15.670 & 0.810 & 19.10 \\
left-256 & 19.456 & 21.660 & 0.822 & 24.27 \\
right-1024 & 267.443 & 270.631 & 3.470 & 77.70 \\
left-1024 & 334.841 & 347.082 & 3.489 & 93.84 \\
right-4096 & 4300.024 & 4480.186 & 14.142 & 305.29 \\
left-4096 & 5953.078 & 5419.346 & 14.304 & 380.36 \\
small-rank-stall & 0.093 & 0.093 & 0.094 & 1.00 \\
raw-recognition-16 & 0.147 & 0.150 & 0.143 & 1.01 \\
raw-recognition-128 & 4.080 & 4.187 & 4.105 & 0.99 \\
raw-recognition-512 & 61.126 & 61.666 & 1.857 & 33.97 \\
raw-recognition-2048 & 1042.500 & 1051.843 & 7.582 & 138.34 \\
pd-recognition-16 & 0.187 & 0.197 & 0.188 & 1.02 \\
pd-recognition-128 & 1.314 & 1.296 & 1.292 & 1.02 \\
pd-recognition-512 & 5.473 & 5.559 & 5.567 & 1.01 \\
\bottomrule
\end{tabular}
\end{center}
Median milliseconds; ratios are medians of paired before/adaptive times.


The ordinary PD examples invoke the linked kernel zero times: earlier
simplification already settles those stabilized unknots. Their paired ratios
are 1.005--1.021 and give no meaningful complete-pipeline improvement.
A/A ratios range from 0.940 to 1.194 across this short experiment; small
apparent gains in unchanged paths fall within that variability. In particular,
the 128-strand right chain has only 127 letters and deliberately stays on the
old branch, while the slightly longer left chain crosses the threshold.
These results support a practical large-chain optimization, not a universal
recognition speedup or a quasi-polynomial general algorithm.
